%%%% CAR-bench IJCAI-ECAI 2026 — Track 2 (Cerebras Fast-Reasoning) technical report.
%%%% Compile from a directory containing ijcai26.sty, named.bst, and track2_report.bib
%%%% (copy them from ../FormattingGuidelines-IJCAI-ECAI-26/, or build there).
%%%% 4-page content limit (competition rule). Keep \linenumbers ON for submission.
%%%% Figure: figures/coroutine_bridge.pdf (repo-root figures/, resolved via \graphicspath).

\documentclass{article}
\pdfpagewidth=8.5in
\pdfpageheight=11in

\usepackage{ijcai26}
\usepackage{times}
\usepackage{soul}
\usepackage{url}
\usepackage[hidelinks]{hyperref}
\usepackage[utf8]{inputenc}
\usepackage[small]{caption}
\usepackage{graphicx}
\usepackage{tikz}
\usetikzlibrary{arrows.meta}
\usepackage{amsmath}
\usepackage{booktabs}
\usepackage[switch]{lineno}

% Comment out for camera-ready.
\linenumbers

\urlstyle{same}

\pdfinfo{
/TemplateVersion (IJCAI.2026.0)
}

\title{Policy as Code: A Coroutine-Bridge Harness for\\ Fast-Reasoning Reliability on CAR-bench}

% Non-anonymous submission. Keep \linenumbers ON for the review submission;
% remove it only for a camera-ready version.
\author{
    Ivan Matveev
    \affiliations
    Proxima Ultra
    \emails
    i.a.matveev@gmail.com
}

\begin{document}

\maketitle

\begin{abstract}
CAR-bench evaluates whether tool-using agents stay reliable under real-world
uncertainty, executing every tool inside the evaluator so that each tool
exchange is a separate agent round-trip. A conventional next-action agent
therefore spends one model call per tool step. We present a coroutine-bridge
harness in which the model's only action is to emit a Python program that
blocks and resumes \emph{in place} across evaluator tool exchanges. This
decouples model invocation from tool round-trips: on the public test split the
agent uses a median of two model calls against seven agent turns per task,
resolving a full multi-turn task in a median of 1.8\,s of model latency on
Cerebras \texttt{gpt-oss-120b}. Because the action surface is executable code,
deterministic CAR-bench policies are encoded directly as logic in the tool
layer rather than as prompt rules, enforcing compliance at zero reasoning cost.
The agent attains Pass\textsuperscript{1} 93.3\% and Pass\textsuperscript{3}
87.3\% on the public test split while remaining well within the Track~2
five-call and token budgets.
\end{abstract}

\section{Introduction}

CAR-bench measures the \emph{consistency} and \emph{limit-awareness} of LLM
agents in an in-car assistant domain of 58 interconnected tools, 19 domain
policies, and 254 public tasks spanning base, hallucination, and
disambiguation categories~\cite{kirmayr2026carbench}. Reliability, not
single-shot capability, is the target: the headline metric is
Pass\textsuperscript{3}, which credits a task only when it passes all three
independent trials.

Track~2 restricts the model to Cerebras-hosted \texttt{gpt-oss} and rewards
\emph{compute-aware} harnessing: architectures that convert fast inference into
low end-to-end latency and reliable behavior. A defining property of the
benchmark shapes what ``compute-aware'' means here. Unlike a standard
Agent-to-Agent (A2A) deployment, where an agent executes its own tools
privately and one request spans the whole task~\cite{a2a2025}, CAR-bench
\emph{externalizes} tool execution to the evaluator so it can score the tool
trajectory and inject environment perturbations. Every batch of tool calls the
agent emits returns as a fresh A2A message carrying tool results. A
conventional ``next-action'' agent, which the reference templates describe as
one model call per assistant step, therefore pays one model call for every
sequential tool step.

Our contribution is a harness that removes this tax. The model emits a single
executable Python program~\cite{wang2024codeact}; when that program calls a
tool wrapper, a coroutine bridge suspends the Python thread, emits the official
tool call, and resumes the \emph{same} stack frame when results arrive ---
without a new model call. One program can thus drive many sequential and
parallel tool round-trips per model invocation. Two consequences follow, and
both are the substance of this report: (i)~model calls, and hence model latency
on the critical path, drop sharply relative to agent turns; and (ii)~because the
action surface is code, deterministic policy is expressed as ordinary logic in
the tools the code calls, so policy compliance costs the model no attention and
no extra reasoning turns.

\section{The Coroutine Bridge}

\subsection{Externalized Tools Force Round-Trips}

In CAR-bench the participant agent never executes vehicle, navigation, weather,
or productivity tools. It returns a data part of \texttt{tool\_calls}; the
evaluator executes them, advances environment state, and returns a new A2A
message of \texttt{tool\_results}. This is what enables trajectory scoring and
hallucination injection, and it is why the benchmark cannot be treated as a
private single-request task. The cost is structural: an $n$-step dependent tool
chain is $n$ separate agent round-trips.

\subsection{Decoupling Model Calls from Tool Round-Trips}

Each A2A \texttt{context\_id} owns a long-lived worker holding one persistent
Python interpreter, a tool bridge, a scratchpad, and a compact transcript. The
model's sole action is \texttt{execute\_python}. Execution proceeds as follows:

\begin{enumerate}
\itemsep2pt
\item The model emits one Python program (one model call).
\item The program runs in the worker thread until it calls a tool wrapper such
  as \texttt{get\_current\_navigation\_state(...)}.
\item The bridge places the official tool call on the outbox and \emph{blocks}
  the Python thread.
\item The evaluator executes the tool and returns a tool-results message.
\item The bridge delivers those results directly into the blocked call, and the
  same stack frame resumes --- \emph{no model call}.
\item The program branches, calls further tools, or emits the user-facing
  response.
\end{enumerate}

\begin{figure*}[t]
\centering
\resizebox{\textwidth}{!}{%
\begin{tikzpicture}[
  font=\small,
  msg/.style={-{Stealth[length=2mm]}, thick},
  ret/.style={-{Stealth[length=2mm]}, thick, dashed},
  call/.style={-{Stealth[length=2mm]}, very thick, blue!70!black},
  callret/.style={-{Stealth[length=2mm]}, very thick, dashed, blue!70!black},
  life/.style={dashed, gray!60},
  hdr/.style={draw, rounded corners, fill=gray!12, minimum height=6mm,
              inner sep=4pt, font=\small\bfseries},
  note/.style={draw, fill=yellow!22, rounded corners, font=\footnotesize,
               inner sep=2.5pt},
  lbl/.style={font=\footnotesize, above, midway},
  step/.style={font=\footnotesize\itshape, fill=blue!6},
]
\def\xE{0}\def\xW{4.4}\def\xM{8.6}\def\xP{13.4}
\node[hdr] (E) at (\xE,0) {Evaluator};
\node[hdr] (W) at (\xW,0) {Worker};
\node[hdr] (M) at (\xM,0) {Model (gpt-oss-120b)};
\node[hdr] (P) at (\xP,0) {Python frame $+$ ToolBridge};
\foreach \x in {\xE,\xW,\xM,\xP} \draw[life] (\x,-0.45) -- (\x,-9.8);

\draw[msg] (\xE,-0.9) -- node[lbl]{user message (A2A)} (\xW,-0.9);
\node[step] at ({(\xW+\xM)/2},-1.35) {internal step 1 of $\leq$5};
\draw[call]    (\xW,-1.9) -- node[lbl]{execute\_python} (\xM,-1.9);
\draw[callret] (\xM,-2.5) -- node[lbl]{Python program} (\xW,-2.5);
\draw[msg] (\xW,-3.1) -- node[lbl]{run program} (\xP,-3.1);
\draw[msg] (\xP,-3.9) -- node[lbl]{batch tool\_calls (A2A turn)} (\xE,-3.9);
\draw[ret] (\xE,-4.5) -- node[lbl]{tool\_results (A2A turn)} (\xP,-4.5);
\node[note, anchor=east] at (\xP-0.25,-5.0) {resume --- no model call};
\draw[msg] (\xP,-5.6) -- node[lbl]{dependent tool\_call (A2A turn)} (\xE,-5.6);
\draw[ret] (\xE,-6.2) -- node[lbl]{tool\_results (A2A turn)} (\xP,-6.2);
\node[note, anchor=east] at (\xP-0.25,-6.7) {resume --- no model call};
\draw[ret] (\xP,-7.3) -- node[lbl]{observation (program done)} (\xW,-7.3);
\node[step] at ({(\xW+\xM)/2},-7.75) {internal step 2 of $\leq$5};
\draw[call]    (\xW,-8.3) -- node[lbl]{execute\_python} (\xM,-8.3);
\draw[callret] (\xM,-8.9) -- node[lbl]{program $\to$ respond()} (\xW,-8.9);
\draw[msg] (\xW,-9.5) -- node[lbl]{final text $+$ turn\_metrics (A2A)} (\xE,-9.5);
\end{tikzpicture}}
\caption{One user turn. Blue = model call; black solid = request, dashed =
return. Two model calls (\texttt{execute\_python}) drive several evaluator tool
round-trips: when a tool result returns, the blocked Python frame resumes
\emph{without} a model call. The internal loop is bounded at five sequential
model calls, matching the Track~2 budget.}
\label{fig:bridge}
\end{figure*}

The model is invoked once per \emph{decision}, not once per tool result.
Figure~\ref{fig:bridge} traces one user turn.
Crucially, this hides nothing: every tool call still appears in the official
A2A trajectory exactly as the evaluator expects, so the technique is fully
compliant. What changes is only the number of \emph{model} invocations between
tool round-trips. Table~\ref{tab:compute} quantifies the effect: a median of
two model calls carries a median of seven agent turns per task.

\subsection{Parallel Batching}

Independent reads are issued together via \texttt{batch([...])}, which emits one
parallel A2A tool-call message and returns results in input order. Dependent
calls remain ordinary sequential Python. Batching collapses independent tool
round-trips in addition to the model-call savings, further shortening the
critical path without sacrificing the natural imperative control flow that
makes code actions expressive~\cite{wang2024codeact}.

\section{Policy as Code}

Many CAR-bench policies are \emph{deterministic given grounded facts}: whether
opening a window above 25\% requires an AC confirmation, whether high beams are
blocked while fog lights are on, which route default applies to a newly created
segment. The CodeAct action surface lets us encode such policy directly as
logic in the tool layer rather than as prose the model must recall. This is the
central design principle of the harness.

We separate two responsibilities. The \emph{runtime} owns facts and strict
policy; the \emph{model} owns interpretation of the user's intent. Concretely,
policy-critical raw wrappers delegate to guarded implementations
(\texttt{set\_fog\_lights} to a fog-light safety helper,
\texttt{open\_close\_window} to an AC-aware helper), so the policy executes
whether the model calls the public name or the helper. A confirmation state
machine stores the fully grounded pending action and resumes the exact call on
the user's next turn. Result normalization gives the model stable field names
so a simple, direct tool path is not silently trapped by a dynamic result key.
Missing capabilities are resolved reactively against the live per-task tool
surface --- a membership test the organizers confirmed as permitted --- never
by comparing against a bundled catalog.

The practical test for each policy is: \emph{is it deterministic given grounded
facts?} If yes, it becomes enforced code in a tool; if it requires genuine
interpretation of the request, it stays in the prompt. A policy expressed only
in the prompt can be forgotten under attention pressure; a policy expressed as
tool code is enforced by construction and, on the critical path, free.

\section{Reliability Mechanisms}

Three further mechanisms address the specific CAR-bench failure classes.
\textbf{Response obligations}: policy-mandated disclosures (e.g.\ a toll-road
mention, a $>$3\textdegree C cross-zone temperature warning) are registered as
obligations that the response step appends only if the model's own text omitted
them, preserving compound-request flexibility without making warnings optional.
\textbf{Unknown-value sentinels}: hallucination tasks can blank fields in
otherwise successful results; the string \texttt{"unknown"} is normalized to a
sentinel that is storage- and serialization-safe but aborts with a
missing-response acknowledgement the moment model code tries to use it in a
meaningful operation, converting a silent fabrication into an honest limitation.
\textbf{Turn-scoped outcome tracking}: mutation failures are keyed by tool and
target arguments and survive every Python block in a user turn, so an optimistic
success claim cannot outrun a failed side effect.

\section{Results}

We evaluate on the public \emph{test} split (125 tasks; 50 base, 50
hallucination, 25 disambiguation) at three trials (375 evaluations), with
Cerebras \texttt{gpt-oss-120b} at temperature~0 and a JSON code-action contract.
The evaluator user simulator and policy judge are Gemini~2.5~Flash, the
organizers' default.

\begin{table}[t]
\centering
\begin{tabular}{lrrr}
\toprule
Per task (test, 3 trials) & Median & Mean & p90 \\
\midrule
Model calls          & 2      & 2.4   & 4     \\
Agent (A2A) turns    & 7      & 8.0   & 12    \\
Model latency (s)    & 1.8    & 2.4   & 4.4   \\
Output tokens        & 1{,}443 & 1{,}850 & 3{,}625 \\
Input tokens         & 78.2k  & 87.8k & 165.6k \\
\bottomrule
\end{tabular}
\caption{Compute and latency profile. Two model calls carry seven agent turns:
the coroutine bridge decouples model invocation from evaluator tool
round-trips. Input token budget is 500k/task; observed mean is 88k.}
\label{tab:compute}
\end{table}

\begin{table}[t]
\centering
\begin{tabular}{lrr}
\toprule
Split & Pass\textsuperscript{1} & Pass\textsuperscript{3} \\
\midrule
Base            & 90.0 & 78.0 \\
Hallucination   & 98.7 & 96.0 \\
Disambiguation  & 90.7 & 88.0 \\
\midrule
Overall         & 93.3 & 87.3 \\
\bottomrule
\end{tabular}
\caption{Reliability on the public test split (\%). Pass\textsuperscript{3}
credits a task only when all three trials pass.}
\label{tab:reliability}
\end{table}

\paragraph{Latency.}
A complete multi-turn task resolves in a median of 1.8\,s of model latency
(Table~\ref{tab:compute}). We report \emph{model} latency deliberately: of the
5{,}351\,s wall time for the full 375-evaluation run, only 899\,s was
agent LLM time; the remainder is the evaluator's own simulator and judge, which
are not attributable to the agent. Fast per-call inference on Cerebras is thus
compounded by a harness that issues few calls, rather than spent on many.

\paragraph{Reliability.}
The agent reaches Pass\textsuperscript{1} 93.3\% and Pass\textsuperscript{3}
87.3\% (Table~\ref{tab:reliability}). The gap between them is
model-compliance variance under a temperature-0 model sensitive to simulator
phrasing, concentrated in the base split; 87\% of tasks pass all three trials.

\paragraph{Compute compliance.}
The internal loop is capped at five sequential model calls per step, matching
the Track~2 budget; observed usage (median two calls, p90 four) stays well
inside it. Mean input+output is under 90k tokens against the 500k budget.
Token and latency accounting is reported through \texttt{turn\_metrics} on each
concluding response; because the accumulator sums across intermediate tool
exchanges and the evaluator records metrics only on non-tool-call responses,
per-task totals are complete.

\paragraph{Development discipline.}
Two rules kept the harness from overfitting the weaker model. First, a change
was accepted only if a stronger model already outperformed a weaker one on the
bare prompt, so gains reflect structure rather than scaffolding. Second, a
persistent \texttt{gpt-oss-120b} failure was first reproduced on a stronger
model: if the stronger model passed, the ceiling was reachable and the fix was
to encode the missing determinism as tool code (Policy as Code); if both
failed, the defect lay in existing harness and was repaired there rather than
patched over. Candidate changes were judged on consistent-failure sets across
runs, never single noisy trials.

\section{Limitations and Conclusion}

The residual Pass\textsuperscript{3} gap is dominated by simulator and judge
variance and by rare paths where a temperature-0 120B model chooses a valid-
looking but wrong action despite enforced structure; a stronger model closes
most of these without additional rules, indicating a model-compliance rather
than architectural ceiling. Worker state is in memory and not recoverable
across process restarts.

The coroutine bridge shows that the most effective way to exploit fast
inference under CAR-bench's externalized-tool design is not more model calls
but a harness that lets each fast call accomplish an entire branching program,
with deterministic policy compiled into the tools that program invokes. The
result is low whole-task latency and reliable, compliant behavior on a
benchmark built to expose the opposite.

\bibliographystyle{named}
\bibliography{track2_report}

\end{document}
